\section{Connecting an agent}
\label{sec:agent-integration}

The endpoint of Section~\ref{sec:remote-access} and the retrieval tools of
Section~\ref{sec:corpus-operation} are both consumed by the same client: an
agent that reasons, calls tools, and reads our documents. This section records
how that client is configured, using Hermes Agent as the worked example. The
principles carry to any agent framework; the key names do not.

\subsection{Three things an agent needs}

\begin{enumerate}
  \item A model that declares \textit{tools} --- see
        Section~\ref{sec:capabilities}. Without it the agent cannot act at all.
  \item A provider entry that actually resolves to our endpoint.
  \item The MCP server, for retrieval over our own documents.
\end{enumerate}

\subsection{The provider must be named}

A custom OpenAI-compatible endpoint is declared as a provider. In Hermes Agent
a \textit{bare} \cmd{provider: custom} does not work: the resolver returns
\cmd{None} for the literal string \cmd{"custom"},

\begin{lstlisting}
if not requested_norm or requested_norm == "custom": return None
\end{lstlisting}

and resolution falls through to the tool's default cloud provider, which is
then sent our local model name together with whatever key is on file. The
symptom is a complaint about a provider nobody configured --- a rejected API
key for a service that was never meant to be used --- and nothing in the error
indicates that a different endpoint was contacted. The remedy is a
\textbf{named} provider entry plus \cmd{custom:<name>}, which gives resolution a
real base URL to find.

\begin{lstlisting}
providers:
  local:
    base_url: http://localhost:11435/v1     # the tunnel, not 11434
    api_key: no-key-required                # Ollama ignores it; must be non-empty
    default_model: "qwen3.5:122b-a10b"

model:
  default: "qwen3.5:122b-a10b"
  provider: custom:local                    # NOT a bare "custom"
  base_url: http://localhost:11435/v1
  context_length: 131072                    # see Section 5.5 on VRAM headroom
  max_tokens: 8192
  timeout_seconds: 300
\end{lstlisting}

Three of those values are worth justifying.

\textbf{Port 11435, not 11434.} On a laptop already running its own Ollama, the
forward cannot bind 11434 and the client silently talks to the local instance
instead --- appearing to work, while using an entirely different model.

\textbf{\cmd{context\_length: 131072}, not the model's native 262\,144.} The KV
cache does not fit the remaining VRAM on a single card
(Section~\ref{sec:measured}).

\textbf{\cmd{max\_tokens: 8192}.} A reasoning model spends tokens before
emitting anything visible; a tight budget returns empty content, which is
indistinguishable from a hung request. A tool-calling turn that also writes an
answer from retrieved passages needs the headroom.

\subsection{Auxiliary slots}

Where the main model lacks a capability, an agent may route that task to a
separate model. Hermes calls these \textit{auxiliary} slots --- vision, web
summarisation and compression among them --- and, importantly, calls them
\textbf{without} a tools array. That makes an auxiliary slot the correct home
for a vision-only model, which would otherwise be rejected outright
(Section~\ref{sec:capabilities}).

\begin{lstlisting}
auxiliary:
  vision:
    model: ""            # empty + "auto" = use the main chat model
    provider: "auto"
    base_url: null
\end{lstlisting}

With \cmd{qwen3.5:122b-a10b} as the main model the slot is left on
\cmd{auto}, because that model has vision of its own. The slot earns its keep
only when the main model does not --- pairing \cmd{gpt-oss:120b}, which has
tools but no vision, with a vision model that has no tools. Each covers the
other's gap, and neither is asked to do what it cannot.

\subsection{Wiring in the corpus}

The MCP server is declared separately from the model, and is reached
\textbf{directly over the office LAN} rather than through the SSH tunnel --- it
is a different service on a different port.

\begin{lstlisting}
mcp_servers:
  corpus:
    url: http://192.168.1.57:8765/mcp
    headers:
      Authorization: "Bearer ${env:CORPUS_MCP_TOKEN}"
    timeout: 300          # rerank with the 120B takes ~20 s; 300 is the default
    connect_timeout: 60
    enabled: true
    tools:
      include: [search_corpus, ask_corpus, corpus_status]
      exclude: []
      resources: false
      prompts: false
\end{lstlisting}

The transport key is omitted deliberately: it is inferred from the presence of
\cmd{url}, and is needed only to force SSE instead of Streamable HTTP. The
token belongs in the agent's environment file, not here, so that the
configuration itself stays safe to copy and to commit.

\subsection{Verifying the wiring}

Check the two halves independently before blaming the agent. The endpoint:

\begin{lstlisting}
curl http://localhost:11435/api/tags          # must list the expected model
\end{lstlisting}

and the MCP server, which needs a session before it will answer --- a bare
\cmd{tools/list} returns \cmd{-32600 Missing session ID}:

\begin{lstlisting}
TOK=$(ssh eugen@192.168.1.57 'cat ~/.config/corpus/mcp-token')
H=(-H "Authorization: Bearer $TOK" -H "Content-Type: application/json"
   -H "Accept: application/json, text/event-stream")

SID=$(curl -s -D- -X POST http://192.168.1.57:8765/mcp "${H[@]}" \
  -d '{"jsonrpc":"2.0","id":1,"method":"initialize","params":{
       "protocolVersion":"2025-06-18","capabilities":{},
       "clientInfo":{"name":"healthcheck","version":"1"}}}' \
  -o /dev/null | tr -d '\r' | awk -F': ' '/^mcp-session-id/{print $2}')

curl -s -X POST http://192.168.1.57:8765/mcp "${H[@]}" \
  -H "mcp-session-id: $SID" \
  -d '{"jsonrpc":"2.0","method":"notifications/initialized"}' -o /dev/null

curl -s -X POST http://192.168.1.57:8765/mcp "${H[@]}" \
  -H "mcp-session-id: $SID" \
  -d '{"jsonrpc":"2.0","id":2,"method":"tools/list"}'
\end{lstlisting}

A request without the bearer token returns HTTP 401, which is the intended
behaviour and a quick confirmation that authentication is enforced.
